% --- CORRECCIÓN IMPORTANTE EN LA PRIMERA LÍNEA ---
% Agregamos "xcolor=table" para que funcione \rowcolor sin errores
\documentclass[aspectratio=169, 10pt, xcolor=table]{beamer}

% --- Configuración del Tema (estilo profesional como el ejemplo) ---
\usetheme{Madrid}
\usecolortheme{dolphin}

% --- Paquetes necesarios ---
\usepackage[utf8]{inputenc}
\usepackage[spanish]{babel}
\usepackage{amsmath, amssymb}
\usepackage{booktabs}
\usepackage{graphicx}
\usepackage{listings}
\usepackage{tikz}
\usetikzlibrary{arrows.meta, positioning, shapes.geometric}

% --- Ruta de Imágenes (crea la carpeta "imagenes" y coloca ahí tus archivos) ---
\graphicspath{{./imagenes/}}

% --- Estilo para código Python ---
\definecolor{codegreen}{rgb}{0,0.6,0}
\definecolor{codegray}{rgb}{0.5,0.5,0.5}
\definecolor{codepurple}{rgb}{0.58,0,0.82}
\definecolor{backcolour}{rgb}{0.95,0.95,0.92}
\lstset{
    language=Python,
    backgroundcolor=\color{backcolour},
    commentstyle=\color{codegreen},
    keywordstyle=\color{blue},
    numberstyle=\tiny\color{codegray},
    stringstyle=\color{codepurple},
    basicstyle=\ttfamily\scriptsize,
    breaklines=true,
    frame=single,
    showstringspaces=false
}

% --- Pie de página personalizado ---
\makeatletter

\setbeamertemplate{footline}{
  \leavevmode%
  \hbox{%
  \begin{beamercolorbox}[wd=.333333\paperwidth,ht=2.25ex,dp=1ex,center]{author in head/foot}%
    \usebeamerfont{author in head/foot}\insertshortauthor
  \end{beamercolorbox}%
  \begin{beamercolorbox}[wd=.333333\paperwidth,ht=2.25ex,dp=1ex,center]{title in head/foot}%
    \usebeamerfont{title in head/foot}Sesión 4
  \end{beamercolorbox}%
  \begin{beamercolorbox}[wd=.333333\paperwidth,ht=2.25ex,dp=1ex,right]{date in head/foot}%
    \usebeamerfont{date in head/foot}NLP \hfill \insertframenumber/\inserttotalframenumber \hspace*{0.3cm}
  \end{beamercolorbox}}%
  \vskip0pt%
}

\makeatother

% --- Datos de la portada ---
\title[SESIÓN 4 NLP]{\textbf{Sesión 4}}
\subtitle{Word Embeddings: Word2Vec y GloVe}
\author[Prof. C. Sánchez]{Carlos César Sánchez Coronel}
\institute{\textbf{NLP}}
\date{}

% Logo opcional (descomentar si tienes la imagen)
% \titlegraphic{\includegraphics[height=1.5cm]{logo.png}}

% --- Eliminar la barra de navegación (opcional) ---
\setbeamertemplate{navigation symbols}{}

% --- Índice al inicio de cada sección ---
\AtBeginSection[]{
  \begin{frame}
    \frametitle{Contenido}
    \tableofcontents[currentsection]
  \end{frame}
}

% ============================================================================
% INICIO DEL DOCUMENTO
% ============================================================================
\begin{document}

% --- Portada ---
\begin{frame}
    \titlepage
\end{frame}

% --- Índice general ---
\begin{frame}{Contenido}
    \tableofcontents
\end{frame}

% ============================================================================
% SECCIÓN 1: LIMITACIONES DE REPRESENTACIONES DISPERSAS
% ============================================================================
\section{Limitaciones de Representaciones Dispersas}

\begin{frame}{Problemas de BOW y TF-IDF}
    \begin{itemize}
        \item \textbf{Alta dimensionalidad}: vocabulario de miles/millones de términos.
        \item \textbf{Esparsidad}: la mayoría de entradas son cero.
        \item \textbf{Falta de semántica}: ``gato'' y ``felino'' son independientes.
        \item \textbf{Pérdida de relaciones}: no capturan analogías (rey - hombre + mujer $\approx$ reina).
        \item \textbf{LSA} (Latent Semantic Analysis) intenta reducir dimensionalidad, pero es lineal y limitado.
    \end{itemize}
    % Basado en PDF y Pregunta 1
\end{frame}

% ============================================================================
% SECCIÓN 2: CONCEPTO DE WORD EMBEDDINGS
% ============================================================================
\section{Word Embeddings}

\begin{frame}{Definición}
    \begin{block}{Word Embedding}
        Representación densa de palabras en un espacio vectorial continuo de baja dimensión (típicamente 50-300).
    \end{block}
    \begin{itemize}
        \item \textbf{Hipótesis distribucional} (Harris 1954): ``palabras que ocurren en contextos similares tienen significados similares''.
        \item Las relaciones semánticas se reflejan como operaciones vectoriales.
    \end{itemize}
    % Basado en PDF y Pregunta 1
\end{frame}

% ============================================================================
% SECCIÓN 3: WORD2VEC
% ============================================================================
\section{Word2Vec}

\begin{frame}{Arquitecturas de Word2Vec}
    \begin{columns}[t]
        \column{0.5\textwidth}
        \textbf{CBOW (Continuous Bag of Words)}
        \begin{itemize}
            \item Predice la palabra objetivo a partir del contexto.
            \item Más rápido, mejor para palabras frecuentes.
        \end{itemize}
        \column{0.5\textwidth}
        \textbf{Skip-gram}
        \begin{itemize}
            \item Predice el contexto a partir de la palabra objetivo.
            \item Mejor para palabras poco frecuentes.
        \end{itemize}
    \end{columns}
    % Basado en PDF y Pregunta 2
\end{frame}

\begin{frame}{CBOW: formalización}
    \begin{block}{Objetivo}
        Dado un contexto de $C$ palabras $\{w_{c1}, \dots, w_{cC}\}$, maximizar $P(w_t | \text{contexto})$.
    \end{block}
    \begin{enumerate}
        \item Embeddings de contexto: $\mathbf{v}_{c_i} = W[w_{c_i}]$.
        \item Promedio: $\mathbf{h} = \frac{1}{C}\sum \mathbf{v}_{c_i}$.
        \item Puntuaciones: $\mathbf{u} = \mathbf{h}^T W'$.
        \item Softmax: $P(w_t|\text{contexto}) = \frac{\exp(u_t)}{\sum_k \exp(u_k)}$.
    \end{enumerate}
    Pérdida: entropía cruzada negativa $-\log P(w_t|\text{contexto})$.
    % Basado en PDF
\end{frame}

\begin{frame}{Skip-gram y negative sampling}
    \begin{block}{Skip-gram}
        Para una palabra objetivo $w_t$, maximizar $P(w_c | w_t)$ para cada contexto $w_c$.
    \end{block}
    \begin{block}{Negative Sampling (aproximación)}
        Para un par $(w_t, w_c)$:
        \[
        L = -\log \sigma(\mathbf{v}_{w_t}^T \mathbf{u}_{w_c}) - \sum_{i=1}^{k} \log \sigma(-\mathbf{v}_{w_t}^T \mathbf{u}_{w_{n_i}})
        \]
        donde $\sigma$ es sigmoide, $w_{n_i}$ son palabras negativas muestreadas.
    \end{block}
    \begin{itemize}
        \item Evita el costoso softmax sobre todo el vocabulario.
        \item Distribución de ruido: $P_n(w) \propto \text{freq}(w)^{3/4}$.
    \end{itemize}
    % Basado en PDF y Pregunta 2, 8
\end{frame}

\begin{frame}{Comparación CBOW vs Skip-gram}
    \begin{table}
        \centering
        \begin{tabular}{l|c|c}
            \toprule
            \textbf{Característica} & \textbf{CBOW} & \textbf{Skip-gram} \\
            \midrule
            Dirección & Contexto $\rightarrow$ Objetivo & Objetivo $\rightarrow$ Contexto \\
            Velocidad & Más rápida & Más lenta \\
            Palabras frecuentes & Bueno & Bueno \\
            Palabras poco frecuentes & Regular & Excelente \\
            Uso típico & Corpus grandes, eficiencia & Corpus pequeños, palabras raras \\
            \bottomrule
        \end{tabular}
    \end{table}
    % Basado en PDF
\end{frame}

% ============================================================================
% SECCIÓN 4: GLOVE
% ============================================================================
\section{GloVe}

\begin{frame}{GloVe: fundamento}
    \begin{block}{Matriz de co-ocurrencia}
        $X_{ij}$ = número de veces que la palabra $j$ aparece en el contexto de $i$ (ventana fija).
    \end{block}
    \begin{block}{Hipótesis}
        \[
        \mathbf{v}_i^T \mathbf{v}_j + b_i + b_j \approx \log X_{ij}
        \]
        donde $\mathbf{v}_i$, $\mathbf{v}_j$ son embeddings y $b_i$, $b_j$ sesgos.
    \end{block}
    % Basado en PDF y Pregunta 12
\end{frame}

\begin{frame}{Función de pérdida de GloVe}
    \[
    \mathcal{L} = \sum_{i,j=1}^{|V|} f(X_{ij}) (\mathbf{v}_i^T \mathbf{v}_j + b_i + b_j - \log X_{ij})^2
    \]
    \begin{block}{Función de ponderación $f$}
        \[
        f(x) = \begin{cases}
            (x/x_{\max})^\alpha & \text{si } x < x_{\max} \\
            1 & \text{en otro caso}
        \end{cases}
        \]
        ($x_{\max}=100$, $\alpha=0.75$) para dar menos peso a co-ocurrencias muy raras o muy frecuentes.
    \end{block}
    % Basado en PDF y Pregunta 12
\end{frame}

\begin{frame}{Ventajas de GloVe}
    \begin{itemize}
        \item Aprovecha estadísticas globales del corpus (matriz de co-ocurrencia).
        \item Buenas propiedades en tareas de analogía.
        \item Entrenamiento eficiente (procesa pares no nulos).
    \end{itemize}
    % Basado en PDF
\end{frame}

% ============================================================================
% SECCIÓN 5: PROPIEDADES DE LOS EMBEDDINGS
% ============================================================================
\section{Propiedades de los Embeddings}

\begin{frame}{Relaciones semánticas y sintácticas}
    \begin{block}{Analogías}
        \[
        \mathbf{v}_{\text{rey}} - \mathbf{v}_{\text{hombre}} + \mathbf{v}_{\text{mujer}} \approx \mathbf{v}_{\text{reina}}
        \]
        Otros ejemplos:
        \begin{itemize}
            \item Madrid - España + Francia $\approx$ París
            \item caminar - camino + nado $\approx$ nadar
        \end{itemize}
    \end{block}
    % Basado en PDF y Pregunta 3, 19
\end{frame}

\begin{frame}{Evaluación intrínseca de embeddings}
    \begin{itemize}
        \item \textbf{Conjuntos de analogías}: Google Analogies Dataset (semánticas y sintácticas).
        \item Se mide la precisión con que el modelo predice la palabra correcta mediante operaciones vectoriales.
        \item Otras evaluaciones: similitud de palabras (WordSim-353), etc.
    \end{itemize}
    % Basado en Pregunta 3
\end{frame}

\begin{frame}{Visualización con t-SNE}
    \begin{itemize}
        \item t-SNE reduce embeddings a 2D preservando la estructura local.
        \item Palabras semánticamente similares aparecen agrupadas.
        \item Precaución: las distancias globales no son interpretables.
    \end{itemize}
    % Basado en PDF y Pregunta 10
    % IMAGEN: Gráfico t-SNE de embeddings
\end{frame}

% ============================================================================
% SECCIÓN 6: EMBEDDINGS CONTEXTUALES
% ============================================================================
\section{Embeddings Contextuales}

\begin{frame}{Estáticos vs Contextuales}
    \begin{columns}[t]
        \column{0.5\textwidth}
        \textbf{Estáticos}
        \begin{itemize}
            \item Un vector por palabra (Word2Vec, GloVe).
            \item No resuelven polisemia (``banco'' = mismo vector).
        \end{itemize}
        \column{0.5\textwidth}
        \textbf{Contextuales}
        \begin{itemize}
            \item Vector depende del contexto (ELMo, BERT, GPT).
            \item Diferentes sentidos tienen diferentes representaciones.
        \end{itemize}
    \end{columns}
    % Basado en PDF y Pregunta 5, 15
\end{frame}

% ============================================================================
% SECCIÓN 7: CASO INTEGRADO - ENTRENAMIENTO CON GENSIM
% ============================================================================
\section{Caso Integrado: Entrenamiento con Gensim}

\begin{frame}[fragile]{Paso 1: Preprocesamiento}
    \begin{lstlisting}
from gensim.models import Word2Vec
from nltk.tokenize import word_tokenize
import nltk
nltk.download('punkt')

sentences = [
    "el gato persigue al ratón",
    "el perro persigue al gato",
    "el ratón huye del gato"
]
tokenized = [word_tokenize(s.lower()) for s in sentences]
    \end{lstlisting}
    % Basado en PDF
\end{frame}

\begin{frame}[fragile]{Paso 2: Entrenar Word2Vec}
    \begin{lstlisting}
model = Word2Vec(
    sentences=tokenized,
    vector_size=100,
    window=5,
    min_count=1,
    workers=4,
    sg=1  # 1 para Skip-gram, 0 para CBOW
)
    \end{lstlisting}
    % Basado en PDF y Pregunta 6
\end{frame}

\begin{frame}[fragile]{Paso 3: Explorar embeddings}
    \begin{lstlisting}
# Embedding de una palabra
vector_gato = model.wv['gato']

# Palabras más similares
model.wv.most_similar('gato', topn=5)

# Analogías: rey - hombre + mujer = ?
result = model.wv.most_similar(positive=['mujer', 'rey'],
                               negative=['hombre'])
    \end{lstlisting}
    % Basado en PDF y Pregunta 3
\end{frame}

\begin{frame}[fragile]{Paso 4: Visualización con t-SNE}
    \begin{lstlisting}
from sklearn.manifold import TSNE
import matplotlib.pyplot as plt

words = list(model.wv.index_to_key)[:100]
vectors = [model.wv[word] for word in words]

tsne = TSNE(n_components=2, random_state=42)
vectors_2d = tsne.fit_transform(vectors)

plt.figure(figsize=(12,8))
for i, word in enumerate(words):
    plt.scatter(vectors_2d[i,0], vectors_2d[i,1])
    plt.annotate(word, (vectors_2d[i,0], vectors_2d[i,1]))
plt.show()
    \end{lstlisting}
\end{frame}

\begin{frame}[fragile]{Paso 5: Cargar GloVe pre-entrenado}
    \begin{lstlisting}
from gensim.scripts.glove2word2vec import glove2word2vec
from gensim.models import KeyedVectors

glove2word2vec('glove.6B.100d.txt', 'glove.6B.100d.word2vec.txt')
glove = KeyedVectors.load_word2vec_format('glove.6B.100d.word2vec.txt')
    \end{lstlisting}
    % Basado en PDF y Pregunta 9
\end{frame}

% ============================================================================
% SECCIÓN 8: PREPARACIÓN PARA ENTREVISTAS
% ============================================================================
\section{Preparación para Entrevistas}

\begin{frame}{Preguntas típicas}
    \begin{itemize}
        \item ¿Qué son los word embeddings y por qué son mejores que BOW/TF-IDF?
        \item Explica CBOW y Skip-gram. ¿Cuándo usarías cada uno?
        \item ¿Qué es negative sampling y por qué es necesario?
        \item ¿Cómo capturan relaciones semánticas los embeddings? (analogías)
        \item Diferencia entre Word2Vec y GloVe.
        \item ¿Qué son los embeddings contextuales? ¿En qué mejoran a los estáticos?
    \end{itemize}
    % Basado en PDF y solucionario
\end{frame}

\begin{frame}{Preguntas avanzadas}
    \begin{itemize}
        \item ¿Cómo funciona la función de pérdida de GloVe?
        \item ¿Qué es OOV y cómo lo manejan FastText o subword embeddings?
        \item ¿Cómo detectar y mitigar sesgos en embeddings?
        \item ¿Cómo evaluarías la calidad de embeddings? (intrínseca vs extrínseca)
    \end{itemize}
    % Basado en Preguntas 7, 11, 12, 16, 17
\end{frame}

% ============================================================================
% SECCIÓN 9: CONEXIÓN CON EL RESTO DEL CURSO
% ============================================================================
\section{Conexión con el resto del curso}

\begin{frame}{Próximos pasos}
    \begin{itemize}
        \item \textbf{Sesión 5}: Clasificación de textos con embeddings.
        \item \textbf{Sesión 6}: Modelos secuenciales (RNN, LSTM) que usan embeddings.
        \item \textbf{Sesión 7}: Transformers (BERT) que generan embeddings contextuales.
    \end{itemize}
\end{frame}

% ============================================================================
% SECCIÓN 10: RESUMEN
% ============================================================================
\section{Resumen}

\begin{frame}{Lo que hemos aprendido}
    \begin{exampleblock}{Conceptos clave}
        \begin{itemize}
            \item \textbf{Limitaciones de representaciones dispersas}: alta dimensionalidad, falta de semántica.
            \item \textbf{Word2Vec}: CBOW (contexto $\rightarrow$ palabra) y Skip-gram (palabra $\rightarrow$ contexto). Negative sampling.
            \item \textbf{GloVe}: factorización de matriz de co-ocurrencia global.
            \item \textbf{Propiedades}: analogías vectoriales, evaluación intrínseca, visualización t-SNE.
            \item \textbf{Embeddings contextuales}: solucionan polisemia (BERT, ELMo).
            \item \textbf{Caso práctico}: entrenamiento y exploración con Gensim.
            \item \textbf{Sesgos}: detección y mitigación.
        \end{itemize}
    \end{exampleblock}
\end{frame}

% --- Diapositiva final ---
\begin{frame}
    \centering
    \Huge \textbf{¡Gracias!}

\end{frame}

\end{document}